\section{BSD as a non-eliminative AOR instance}
\label{sec:aor-instance}

The conditional landing of \Cref{sec:landing} gives the scalar
Strong-BSD identity under the standing shell-closure record, the named
recognition sources, and the headline imports.
This final technical section records the corresponding
Audited Operational Realisability reading.  In the sense of
\cite{TsiokosAOR2026}, an AOR instance is a refinement-stable,
audit-closed membership predicate: informally, the framework's
``holds in reality'' assertion for an object, meaning that the
residual defects left by refinement have all been accounted for by
approved records.  The reading here is non-eliminative.  The open
arithmetic content is not erased; it is recorded as recognition-source
content, while the mechanical pieces and imported theorem stacks are
accounted for by construction or by approved external citation.

\begin{definition}[BSD AOR-instance object]\label{def:closure:aor-sel-instance}
The AOR-instance object attached to \(\SelBSD\) is the carrier whose
membership witness is the bundle of closure records for the BSD trace
shell.  The bundle has two kinds of entries.  The mechanical entries
are the construction of \(\SelBSD\), the \(\kapr\) normalization, and
the imported theorem stacks recorded through an approved
external-record citation chain.  The substantive entries are the three
recognition-source records of \Cref{sec:recognition-sources}.  The
object is built using the AOR primitives of \cite{TsiokosAOR2026}:
refinement-stable membership, structural-defect accounting, and
discharge records.
\end{definition}

\begin{theorem}[Mechanical-record discharges]\label{thm:closure:aor-mechanical-records}
The mechanical membership components discharge cleanly.  The
\(\SelBSD\) construction and the \(\kapr\) normalization are
discharged by construction.  The \(\chiCTp\), \(\AofE\), and
Beilinson import structures are discharged by an approved
external-record citation chain: they are cited as imported theorem
stacks, not reproved.

For the normalization record, we use \citet{TsiokosBSDApparatus2026}.
Fix the cascade convention
\[
  c_E := \Tam(E)=\prod_{\ell} c_{\ell}(E).
\]
Then the cascade leading-term form
\[
  \frac{L^{(r)}(E,1)}{r!}
  =
  \kapr(E)\,\Omega_E\,\Reg_r(E)\,c_E
\]
agrees with the standard Strong-BSD leading-coefficient form
\[
  \frac{L^{(r)}(E,1)}{r!}
  =
  \Omega_E\,\Reg_r(E)\,
  \frac{\#\Sha(E/\Q)\,\Tam(E)}{\#E(\Q)_{\tors}^{\,2}}
\]
if and only if
\[
  \kapr(E)=
  \frac{\#\Sha(E/\Q)}{\#E(\Q)_{\tors}^{\,2}}.
\]
This is a restatement in cascade coordinates, not a proof of the
Strong-BSD identity.
\footnote{Verified in Lean as
\texttt{SixBirdsBSD.Closure.AORInstance.aorMechanicalRecords};
see App~\ref{app:formalization}.}
\end{theorem}

\begin{proof}
The shell itself is one of the constructed objects of the closure
architecture, so its mechanical entry is accounted for by construction.
The displayed \(\kapr\) equivalence is likewise a constructed
normalization record: once the cascade convention \(c_E=\Tam(E)\) is
fixed, it identifies the scalar that makes the cascade leading-term
form and the standard Strong-BSD leading-coefficient form the same
formula.  It does not assert that the formula is true independently of
the conditional landing.

It remains only to account for the imported theorem stacks.  The
\(\chiCTp\) Cassels--Tate stack, the \(\AofE\) Selmer-complex pairing
stack, and the Beilinson comparison stack enter as named external
records.  They are therefore discharged by an approved citation chain,
not by a derivation inside the closure.  These three citation records,
together with the constructed shell and the \(\kapr\) normalization,
exhaust the mechanical part of the AOR membership witness.
\end{proof}

\begin{theorem}[Recognition-source discharges]\label{thm:closure:aor-recognition-discharge}
The three recognition sources discharge as bridged records.  Here a
bridged record means a recognition residual that is reclassified as
supplied closure content rather than eliminated by a derivation.  Each
\(\Gamma_{\mathrm{BSD}}^{*}\) is a primary structural defect of
source type, carrying its forced secondary defects of role, target,
transport, and limit, and is discharged by such a bridged record.  In
particular, none of the three recognition sources is eliminated.
\footnote{Verified in Lean as
\texttt{SixBirdsBSD.Closure.AORInstance.aorRecognitionDischarge};
see App~\ref{app:formalization}.}
\end{theorem}

\begin{proof}
Each of the three recognition sources is an open arithmetic obligation
that the closure names and consumes but does not derive.  In the audit
presentation, such an obligation is therefore recorded as a
source-type structural defect, together with the secondary information
that specifies its role in the argument, its target factor in the BSD
residual, its transport through the shell, and the limit of what the
closure claims about it.

The discharge is by bridging.  That is, the record says that the
obligation is supplied as recognition content and is available to the
conditional landing; it does not say that the obligation has been
proved away.  This applies to the \(p\)-adic descent source, the
signed-Selmer persistence source, and the higher-rank fixity source.
Thus all three recognition records are accounted for, and all three
remain non-eliminated.
\end{proof}

\Cref{tab:aor-discharge-register} lists the mechanical, imported, and
recognition-source atoms that enter the local AOR membership witness.
\begin{table}[tbp]
\centering
\small
\caption{The eight AOR discharge atoms for the BSD trace shell.  The
recognition-source rows are bridged, not zeroed.}
\label{tab:aor-discharge-register}
\begin{tabularx}{\textwidth}{@{}>{\raggedright\arraybackslash}p{0.31\textwidth}
>{\raggedright\arraybackslash}p{0.24\textwidth}
>{\raggedright\arraybackslash}X@{}}
\toprule
Residual atom & Primary type & Status \\
\midrule
\(\SelBSD\) construction & mechanical & by construction \\
\(\kapr\) normalization & mechanical & by construction, citing the
apparatus paper \\
\(\chiCTp\) import stack & imported theorem stack & approved external
record \\
\(\AofE\) Selmer-complex substrate & imported theorem stack & approved
external record \\
Beilinson rank-\(\geq2\) comparison & imported theorem stack & approved
external record \\
\(\GammaPD\) & recognition source & bridged, not eliminated \\
\(\GammaSP\) & recognition source & bridged, not eliminated \\
\(\GammaGZ\) & recognition source & bridged, not eliminated \\
\bottomrule
\end{tabularx}
\end{table}


\begin{theorem}[AOR-instance membership]\label{thm:closure:aor-instance}
\(\SelBSD\) is refinement-stable audit-closed, hence an AOR instance,
under the declared presentation: its declared audit records discharge
the structural defects that remain under refinement.  This is the
framework's ``Strong BSD holds in reality'' statement, and it is
non-eliminative.  The recognition records are reclassified as bridged
records, not discharged of their open arithmetic.
\footnote{Verified in Lean as
\texttt{SixBirdsBSD.Closure.AORInstance.aorInstance} over a narrowed
representation -- the local syntactic refinement-stable predicate, not
the full Audited Operational Realisability meta-theory; the narrowing
is recorded in App~\ref{app:formalization}.}
\end{theorem}

\begin{proof}
By \Cref{thm:closure:aor-mechanical-records}, the mechanical
components of the BSD trace shell are accounted for: the shell
construction and the \(\kapr\) normalization are discharged by
construction, while the imported theorem stacks are discharged by
approved external citation.  By
\Cref{thm:closure:aor-recognition-discharge}, the three recognition
sources are accounted for as bridged records.

These two groups exhaust the declared structural defects that persist
under refinement of the BSD trace shell.  The mechanical records are
closed by construction or citation; the recognition records are
carried forward as supplied recognition content.  Therefore every
remaining defect is matched by a declared audit record, so the shell is
refinement-stable audit-closed under the declared presentation.

The last clause is essential.  The membership statement does not
convert the recognition sources into classical proofs of the open
arithmetic obligations.  It records that those obligations are present,
named, and bridged in the audit, which is exactly the non-eliminative
AOR reading.
\end{proof}

\begin{remark}[AOR partial status]\label{rmk:closure:aor-partial-status}
The AOR-instance reading does not close the open obligations.  The
rank-\(\geq 2\) generalized Gross--Zagier identity and the additive
\(p\leq 3\) Iwasawa main conjecture remain open in the literature.
What the AOR reading records is that the BSD carrier is audit-closed
under the declared refinement-stable presentation, with the
recognition records reclassified as bridged content rather than
discarded.  Only the local syntactic refinement-stable predicate is
mechanized here, not the full AOR meta-theory.
\end{remark}
